% !TeX root = ../technical_report.tex
\chapter{DVC Dataset Versioning and DagsHub Remote Storage}
\label{ch:dvc_dagshub_storage}

The physical topography and bounded execution strategies delineated in Chapter~\ref{ch:bare_metal_host} rigorously address the computational constraints of bare-metal vision architectures. However, ensuring strict provenance and absolute mathematical reproducibility requires an equivalent operational discipline in data management. Modern Machine Learning Operations (MLOps) platforms inevitably confront the impedance mismatch between the versioning mechanics of source code and the storage requirements of large binary assets. Specifically, traditional Git repositories fundamentally degrade when subjected to the multi-gigabyte churn associated with continuous image corpus ingestion and iterative tensor checkpointing. Attempting to force terabytes of heterogeneous biological imagery through the standard Git tree objects leads to catastrophic repository bloat, unacceptable clone times, and severe I/O bottlenecks. This chapter meticulously formalises the integration of Data Version Control (DVC) and DagsHub S3-compatible remote storage as the cryptographic and physical backbone of the TaxonVision-MLOps data topology, ensuring that empirical workflows remain strictly deterministic.

\section{Data Version Control (DVC) Core Architecture}
\label{sec:dvc_core_architecture}

\begin{figure}[htbp]
\centering
\definecolor{DvcTeal}{HTML}{00897B}
\definecolor{DvcAmber}{HTML}{FFA000}
\definecolor{DvcPurple}{HTML}{8E24AA}
\definecolor{DvcSlate}{HTML}{37474F}
\definecolor{DvcBg}{HTML}{E0F2F1}

\resizebox{\textwidth}{!}{%
\begin{tikzpicture}[
    node distance=1.1cm and 1.6cm,
    font=\footnotesize,
    >=stealth,
    gitbox/.style={
        rectangle,
        rounded corners=4pt,
        draw=DvcTeal,
        fill=DvcTeal!8,
        thick,
        inner sep=7pt,
        align=left,
        drop shadow={opacity=0.15, shadow xshift=2pt, shadow yshift=-2pt},
        minimum width=5.2cm
    },
    casbox/.style={
        rectangle,
        rounded corners=4pt,
        draw=DvcAmber,
        fill=DvcAmber!10,
        thick,
        inner sep=7pt,
        align=left,
        drop shadow={opacity=0.15, shadow xshift=2pt, shadow yshift=-2pt},
        minimum width=5.2cm
    },
    remotebox/.style={
        rectangle,
        rounded corners=4pt,
        draw=DvcPurple,
        fill=DvcPurple!8,
        thick,
        inner sep=7pt,
        align=left,
        drop shadow={opacity=0.15, shadow xshift=2pt, shadow yshift=-2pt},
        minimum width=5.2cm
    },
    groupbox/.style={
        rectangle,
        rounded corners=6pt,
        draw=DvcSlate!50,
        fill=DvcBg!30,
        dashed,
        inner sep=10pt
    },
    arrow/.style={
        ->,
        thick,
        draw=DvcSlate
    }
]

% Column 1: Git Tracked State
\node[gitbox] (gitpointer) {
    \textbf{Git Working Tree Pointer File}\\
    \scriptsize \texttt{data/sample.dvc}\\
    \scriptsize $\bullet$ \texttt{md5: f546997db8b542e87e4f006326d04aac.dir}\\
    \scriptsize $\bullet$ \texttt{size: 37744}, \texttt{nfiles: 1}\\
    \scriptsize $\bullet$ \texttt{path: sample}\\
    \scriptsize Lightweight text pointer committed to Git
};

\node[gitbox, below=0.7cm of gitpointer] (lockfile) {
    \textbf{Pipeline DAG Ledger (\texttt{dvc.lock})}\\
    \scriptsize $\bullet$ Cryptographic hash of \texttt{data/manifests}\\
    \scriptsize $\bullet$ Preserved as \texttt{dvc\_dataset\_hash} in MLflow\\
    \scriptsize $\bullet$ Atomic link: Git commit SHA $\leftrightarrow$ Dataset snapshot\\
    \scriptsize $\bullet$ Satisfies MLOps Reproducibility Rule 01
};

% Column 2: Content-Addressed Storage
\node[casbox, right=1.6cm of gitpointer] (localcas) {
    \textbf{Local Content-Addressed Storage}\\
    \scriptsize \texttt{.dvc/cache/}\\
    \scriptsize $\bullet$ 2-character directory sharding:\\
    \scriptsize \phantom{$\bullet$ } \texttt{.dvc/cache/f5/46997db8b542...dir}\\
    \scriptsize $\bullet$ Prevents ext4 inode exhaustion\\
    \scriptsize $\bullet$ Hardlink/Symlink materialisation into workspace
};

\node[remotebox, below=0.7cm of localcas] (dagshub) {
    \textbf{Remote CAS: DagsHub S3 Endpoint}\\
    \scriptsize \texttt{https://dagshub.com/foersben/...s3}\\
    \scriptsize $\bullet$ Bucket: \texttt{s3://dvc}\\
    \scriptsize $\bullet$ Unified token authentication (\texttt{DAGSHUB\_TOKEN})\\
    \scriptsize $\bullet$ Ephemeral agent hydration via \texttt{dvc pull}\\
    \scriptsize $\bullet$ Fail-soft fallback on network partition
};

% Column 3: Workflows
\node[casbox, right=1.6cm of localcas] (workflow) {
    \textbf{Data Lifecycle Protocols}\\
    \scriptsize $\bullet$ \texttt{dvc push}: Uploads novel MD5 chunks only\\
    \scriptsize $\bullet$ \texttt{dvc pull}: Hydrates CI runner workspace\\
    \scriptsize $\bullet$ \texttt{dvc status}: Real-time drift detection\\
    \scriptsize $\bullet$ \texttt{dvc gc}: Prunes unreferenced local blobs
};

% Groupings
\begin{scope}[on background layer]
    \node[groupbox, fit=(gitpointer)(lockfile), label={[font=\bfseries, text=DvcTeal]above:Git Metadata Plane}] (g1) {};
    \node[groupbox, fit=(localcas)(dagshub), label={[font=\bfseries, text=DvcAmber]above:Content-Addressed Storage (CAS)}] (g2) {};
    \node[groupbox, fit=(workflow), label={[font=\bfseries, text=DvcPurple]above:Synchronisation Operations}] (g3) {};
\end{scope}

% Arrows
\draw[arrow] (gitpointer) -- (localcas) node[midway, above, font=\tiny] {Resolves};
\draw[arrow] (lockfile) -- (dagshub) node[midway, below, font=\tiny] {Verifies};
\draw[arrow] (localcas) -- (dagshub) node[midway, right, font=\tiny] {\texttt{push} / \texttt{pull}};
\draw[arrow] (localcas) -- (workflow);

\end{tikzpicture}%
}

\caption{Data Version Control (DVC) content-addressed storage and DagsHub remote synchronization.}
\label{fig:dvc_dataset_versioning}
\end{figure}

At its core, DVC functions by decoupling the versioning of logical pointers (which remain in Git) from the storage of physical binary assets (which are relegated to a dedicated, high-bandwidth backend). This fundamental decoupling ensures that the Git repository remains strictly text-based, compact, and blazingly fast. Meanwhile, vast datasets, ONNX runtime graphs, and Apache Parquet metadata manifests are securely managed via a robust Content-Addressed Storage (CAS) mechanism. This division of responsibility fundamentally mitigates the scaling limitations that typically paralyse monolithic source control paradigms during data-intensive biological modelling.

\subsection{Content-Addressed Storage (CAS) Mechanics}
The absolute cornerstone of DVC is its CAS implementation. Instead of identifying a file by its hierarchical path (which is inherently mutable and context-dependent), CAS identifies every binary artefact strictly by a cryptographic hash of its immutable byte sequence. Upon executing a tracking command (e.g., \texttt{dvc add}), DVC computes the exact MD5 checksum of the target file or directory. The physical data is subsequently ingested into the local cache directory, \texttt{.dvc/cache/}, thereby guaranteeing that mathematically identical data always maps to the same deterministic storage location regardless of its logical nomenclature.

To prevent catastrophic filesystem slowdowns caused by massive directory listings (a known limitation of standard inode structures like \texttt{ext4}), DVC implements a deterministic, two-character subdirectory sharding protocol. For example, if an artefact yields an MD5 hash beginning with \texttt{f5}, its physical cache path becomes \texttt{.dvc/cache/f5/}. The remaining 30 characters of the hash form the filename within that specific shard. This structure explicitly avoids unbounded flat directories and enables instantaneous hash-based retrieval during CI/CD hydration. In the case of directory tracking, the hash corresponds to a structured JSON manifest enumerating the individual file hashes within that directory, suffixed with \texttt{.dir} to logically denote its structural nature.

\subsection{Structural Anatomy of DVC Pipeline Tracking and Pointer Files}
The physical decoupling process generates lightweight, YAML-formatted pointer files (such as atomic \texttt{.dvc} files or the comprehensive \texttt{dvc.yaml} pipeline DAG with its state ledger \texttt{dvc.lock}), which are subsequently committed to Git alongside the source code. While individual \texttt{.dvc} files can track isolated local assets, the full multi-stage Epistemic Diamond DAG is codified within \texttt{dvc.yaml} and locked in \texttt{dvc.lock}, tracking code dependencies, execution parameters, and output artifacts:

\begin{listing}[H]
\begin{minted}{yaml}
schema: '2.0'
stages:
  extract_features:
    cmd: pixi run -e dev python -m taxon_vision.pipelines.extract_features
    deps:
    - data/manifests/dataset_manifest.parquet
    - src/taxon_vision/pipelines/extract_features.py
    params:
      params.yaml:
        extract_features:
          backbone: timm/vit_base_patch16_clip_224.openai
          split_ratio: [0.7, 0.1, 0.1, 0.1]
    outs:
    - md5: 8a4c1f930e1b4279b9087c2b5f001234
      size: 15288320
      path: data/embeddings/train_embeddings.pt
    - md5: 3d7f9b201e5a4891a8761e0b9c345678
      size: 2184000
      path: data/embeddings/val_embeddings.pt
\end{minted}
\caption{Cryptographic DAG and artifact tracking structure within \texttt{dvc.lock}.}
\label{lst:dvc_lock}
\end{listing}

As thoroughly detailed in Listing~\ref{lst:dvc_lock}, the tracking specification intrinsically binds each logical artifact path (\texttt{data/embeddings/train\_embeddings.pt}) to its exact cryptographic MD5 hash, byte size, and stage dependencies. This structure enables Git to strictly track the exact state of manifests, split embeddings, and model checkpoints through time without ever persisting heavy binary payloads in the Git object store. The physical cache shard for each artifact resides within \texttt{.dvc/cache/}, linking declarative pipeline states to immutable disk blocks.

\section{DagsHub S3 Remote Integration}
\label{sec:dagshub_s3_remote}

While the local CAS operates effectively for individual developer workstations and localized prototyping, distributed team collaboration and ephemeral Continuous Integration (CI) pipelines require a centralized, high-availability remote storage layer. To fulfil this absolute requirement comprehensively, the repository natively integrates with DagsHub, establishing a robust, enterprise-grade S3-compatible backend tailored specifically for automated machine learning workloads.

\subsection{Configuration and Protocol Negotiation}
The connection topology is rigidly defined within the hidden \texttt{.dvc/config} file. The configuration instructs the local DVC client to utilise the DagsHub remote as its primary upstream origin, ensuring a unified data fetching strategy across all distributed agents and CI runners.

\begin{listing}[H]
\begin{minted}{ini}
[core]
    remote = origin
['remote "origin"']
    url = s3://dvc
    endpointurl = https://dagshub.com/foersben/taxon-vision-mlops.s3
\end{minted}
\caption{S3 Remote configuration defined in \texttt{.dvc/config}.}
\label{lst:dvc_config}
\end{listing}

As demonstrated in Listing~\ref{lst:dvc_config}, the default remote is explicitly named \texttt{origin}, intentionally mirroring Git semantics for cognitive continuity among developers. The underlying protocol is strictly defined as an S3 bucket (\texttt{s3://dvc}), yet the network traffic is deterministically routed to the custom DagsHub endpoint URL. This approach elegantly circumvents the requirement for expensive, standalone AWS S3 infrastructure while maintaining strict compatibility with established S3 client libraries (such as \texttt{boto3} in Python). This architectural choice drastically reduces operational complexity and financial expenditure for the open-source platform.

\subsection{Authentication and Data Synchronisation Semantics}
Authentication to the S3-compatible endpoint is dynamically negotiated without exposing static AWS-style credentials in plain text. Instead, the pipeline leverages Hash-based Message Authentication Code (HMAC) credentials derived dynamically from a unified access token. In the Jenkins CI environment, the \texttt{DAGSHUB\_TOKEN} credential is simultaneously mapped to both \texttt{AWS\_ACCESS\_KEY\_ID} and \texttt{AWS\_SECRET\_ACCESS\_KEY}. This ingenious duality allows standard S3 libraries to transparently authenticate against the DagsHub proxy, facilitating seamless read and write operations across the entire network cluster.

Data synchronisation is rigorously managed via specific declarative commands, ensuring strict state alignment between local caches and remote storage:
\begin{itemize}
    \item \texttt{dvc push}: Evaluates the delta between the local \texttt{.dvc/cache/} and the remote DagsHub endpoint. Only novel, un-cached binary blobs (identified by missing MD5 hashes on the remote) are transmitted over the network. This intrinsically minimises egress bandwidth and vastly accelerates the iterative model tuning loop.
    \item \texttt{dvc pull}: Executes the reverse operation. It parses the \texttt{.dvc} pointers currently checked out in the Git working tree, queries the remote for the corresponding MD5 hashes, and downloads any missing blobs into the local cache. The blobs are then physically linked (via hardlinks or symlinks, depending on filesystem capabilities) into the working directory, materialising the dataset required for the training process.
    \item \texttt{dvc status}: Provides a real-time differential analysis, instantly reporting discrepancies between the Git-tracked pointers and the physically present local cache or remote storage. This serves as the primary auditing tool for confirming operational integrity before executing pipeline commands.
\end{itemize}

\section{Reproducibility Invariants and Pipeline Provenance}
\label{sec:reproducibility_invariants}

The paramount architectural directive governing this platform is MLOps Reproducibility Rule 01: ensuring exact data-model lineage. Every emitted model checkpoint, particularly the high-value ONNX computational graphs, must be mathematically and cryptographically linkable to a precise, immutable snapshot of the underlying training data. Without this mathematical invariant, debugging catastrophic failure modes, interpreting hyperparameter search spaces, or rectifying epistemic drift in biological classification becomes effectively impossible.

The explicit and automated interplay between DVC and Git inherently enforces this strict invariant. When an automated training pipeline concludes, the resulting weights\\(e.g., \texttt{models/checkpoints/head.pt}) and intermediate states are immediately registered via \texttt{dvc add}, generating a new `.dvc` pointer. This pointer, alongside any updated dataset pointers and training configurations, is committed to Git in a single, atomic transaction. Consequently, checking out a specific Git commit hash instantly and deterministically reconstructs the exact combination of source code, Bayesian hyperparameters, and the precise dataset state utilised for that specific model iteration.

Managing the physical footprint of this provenance requires robust cache pruning protocols. Incremental additions, particularly the partitioning of large Parquet datasets, generate substantial cache churn over time. Because DVC never modifies a file in the CAS (it strictly creates new hashes), obsolete datasets accumulate indefinitely unless actively garbage collected. Routine execution of \texttt{dvc gc --workspace} ensures that only the binary blobs referenced by the currently active Git branch (or explicitly preserved tags) are retained locally. This procedure is absolutely crucial for maintaining operational efficiency, preventing the local NVMe drive from reaching its capacity limits under the storage partitions detailed in Chapter~\ref{ch:gpu_virtualisation}.

\section{CI/CD Integration and Ephemeral Agent Hydration}
\label{sec:ci_cd_integration}

The practical execution of this architecture is prominently demonstrated within the Jenkins CI/CD pipeline. The ephemeral Kubernetes pods (specifically the \texttt{ml-runner} container) initialise in a completely stateless and pristine condition. In order to execute integration tests, compute Class-Balanced Loss metrics (Chapter~\ref{ch:class_balanced_loss}), or rigorously validate the zero-downtime inference deployment, the pod must rapidly hydrate its working directory with the required binary assets before actual execution can commence.

\subsection{Hydration Mechanics in the Jenkinsfile}
The initialisation sequence occurs immediately within the \texttt{stage('Prepare Toolchain')} block of the \texttt{Jenkinsfile}. After the Pixi multi-environment toolchain perfectly constructs the required Python dependency graph, the agent invokes DVC to fetch the necessary datasets and model checkpoints from DagsHub.

\begin{listing}[H]
\begin{minted}{bash}
echo ">>> Pulling DVC tracked datasets and model checkpoints from DagsHub..."
pixi run --frozen -e ci-dev dvc pull || echo "WARNING: DVC pull skipped or failed; continuing with workspace cache."
\end{minted}
\caption{Ephemeral agent hydration logic derived from the \texttt{Jenkinsfile}.}
\label{lst:jenkins_hydration}
\end{listing}

As explicitly illustrated in Listing~\ref{lst:jenkins_hydration}, the invocation utilises the isolated \texttt{ci-dev} environment managed by Pixi. The command interrogates the checked-out \texttt{dvc.lock} manifest and efficiently streams only the required MD5 blobs from the S3 remote, thus flawlessly reconstructing the necessary environment for model testing. Furthermore, within the subsequent \texttt{stage('Model Training \& Checkpoint (GPU)')}, trained outputs are synchronised back to remote storage via \texttt{pixi run --frozen -e dev dvc push models/checkpoints/head.pt models/checkpoints/model\_int8.onnx || true}.

\subsection{Fail-Soft Caching Mechanisms}
A critical requirement for enterprise-grade MLOps is unyielding resilience against external infrastructure degradation. If the DagsHub S3 endpoint experiences sudden latency spikes, degraded throughput, or absolute downtime, the pipeline must not catastrophically stall if the necessary assets are already present on the local cluster. The explicit shell logical OR operator (\texttt{||}) implements a strict fail-soft mechanism. If the network request fails, the pipeline automatically emits a non-fatal warning and continues execution by safely utilising whatever files are currently available in the cached storage layer. This resilient strategy heavily relies on the persistence of the underlying Kubernetes PersistentVolumeClaim (mounted specifically at \texttt{/root/.cache/rattler/cache} and the workspace), allowing successive, ephemeral pod executions to leverage previously downloaded, immutable chunks from prior CI runs. This methodology fundamentally preserves CI/CD throughput even during severe remote storage outages, guaranteeing high availability and unhindered workflow continuation.
